\subsubsection{树链的交}
\inputminted{cpp}{sections/Tree/assets/yzh/path_inter.cpp}

\subsubsection{树的计数 Prufer 序列}
Prufer 编码长度为 ${n-2}$，且度数为 $d_i$ 的点在编码中出现 ${d_i-1}$ 次。
由树得到序列时，每次删除当前编号最小的叶子并记录其相邻点；反向构造时根据编号出现次数恢复各点度数。

$n$ 个点的完全图中，指定各点度数依次为 $d_1,d_2,\ldots,d_n$ 的生成树数量为
\[
  \frac{(n-2)!}{(d_1-1)!(d_2-1)!\cdots(d_n-1)!}.
\]
左右两侧分别有 $n_1,n_2$ 个点的完全二分图，其生成树数量为
$n_1^{n_2-1}n_2^{n_1-1}$。

\subsubsection{有根树计数}
无标号有根树数量满足
\[
  a_{n+1}=\frac1n\sum_{k=1}^{n}\left(\sum_{d\mid k}d\,a_d\right)a_{n-k+1}.
\]

\subsubsection{无根树计数}
$n$ 为奇数时，有 $a_n-\sum_{i=1}^{\lfloor n/2\rfloor}a_i a_{n-i}$ 种不同的无根树；
$n$ 为偶数时，再加上 $\frac12a_{n/2}(a_{n/2}+1)$。
